\section{Discussion}
The results separate two senses of ``representing a label.'' In the category-recognition sense, the member-derived KC vectors work: they retrieve member questions, produce positive coherence margins, and partially align with label-name groupings. In the downstream functional sense, they are uneven. They show some transition-prediction signal, weak difficulty signal, weak graph alignment, and no reliable DKT initialization gain. The same vector can therefore be a reasonable category representation, a useful feature for some auxiliary functionality, and still a poor functional prior for the central KT model.

This distinction matters beyond education. Member aggregation is common because it is simple and intuitive: if a label is defined by its examples, average the examples. But a downstream model may need a relation that is not the same as topical similarity. For KCs, DKT benefits from structure about student trajectories, correctness patterns, skill dependencies, and dataset-specific sequencing. A question about two-step linear equations and a question about simplifying expressions may be textually similar, but their role in a student's learning path can differ. Conversely, two KCs may be functionally related because students transition between them, even when their surface text differs.

The negative result should not be read as a claim that embeddings are useless for education. It says that fixed-category representation by member aggregation is insufficient for this downstream role. A separate line of work, corresponding to Chapter 6, asks whether embeddings can help diagnose or redefine KC boundaries. That is a different claim: instead of representing existing labels as functional priors, it changes the label space. The present paper intentionally keeps that work separate so the Chapter 5 result remains interpretable as evidence about member-derived label representation.

For NLP practice, the implication is methodological. Label representations should be evaluated at two levels: whether they recover the intended label/category, and whether they encode the relation required by each downstream use. The non-DKT diagnostics matter because they prevent an over-simple story: member-derived labels do encode some functional signals. The DKT result matters because the strongest intended use still fails. Passing one downstream diagnostic does not license a broader functional-representation claim.
